\section{Framework Implementation}
\label{sec:implementation}

SandboxBench is about 5{,}100 lines of Go, with a further 700 of tests. Go rather than Rust because the
framework is a distributed load generator and an adapter surface, not a
virtualisation engine: goroutines and contexts express open-loop and
closed-loop generation directly, and the cost of contributing a new provider
adapter matters more to the project than the runner's own peak performance.
Measurements that need precision run \emph{inside} the sandbox
(\S\ref{sec:impl:timing}), where the runner's garbage collector cannot reach
them.

\subsection{Three layers}

\begin{description}
  \item[Provider adapter] owns resource lifecycle and command execution. It
  knows nothing about what is being measured.
  \item[Workload] decides what to measure. It receives either a ready sandbox
  or, when it needs to control lifecycle itself, the provider.
  \item[Runner] decides how to drive it: repetitions, concurrency, arrival
  rate, duration, timeouts, teardown.
\end{description}

The provider contract is five operations on a sandbox --- \texttt{Exec},
\texttt{WriteFile}, \texttt{ReadFile}, \texttt{Suspend}, \texttt{Resume},
\texttt{Destroy} --- plus \texttt{Create} and \texttt{Capabilities} on the
provider. \texttt{Suspend} and \texttt{Resume} return a sentinel
\texttt{ErrUnsupported}, which the runner turns into the \emph{unsupported}
status rather than into a failure or a zero.

Capabilities are checked \emph{before} any sandbox is created, so an
unsupported benchmark costs nothing rather than wasting a create.

\subsection{Three exec transports}
\label{sec:impl:transports}

Implementing three adapters surfaced a finding that we did not expect and that
shapes the rest of the paper: the three providers execute commands over three
different mechanisms, with three different per-command cost structures
(Table~\ref{tab:transports}).

\begin{table}[t]
\caption{Command execution mechanisms. Compiled from each provider's published
API specification while writing the adapters.}
\label{tab:transports}
\small
\begin{tabular}{@{}llp{3.4cm}@{}}
\toprule
Provider & Transport & Per-command cost \\
\midrule
Daytona & REST            & one HTTP request to a per-sandbox toolbox proxy \\
E2B     & Connect-RPC     & a server-streaming RPC to \texttt{envd} inside the sandbox \\
Runta   & SSH             & a session on a multiplexed connection through an SSH gateway \\
\midrule
local   & process         & \texttt{os/exec} on the runner host (reference floor) \\
\bottomrule
\end{tabular}
\end{table}

Runta's public REST API has no execute endpoint at all; it covers lifecycle,
files and checkpoints, and commands go in over the SSH gateway. A suite that
reported a single ``command latency'' column across these three would mostly be
ranking transports. The transport is therefore recorded with every result and
printed in every capability table.

\subsection{Separating the transport from the sandbox}
\label{sec:impl:timing}

Every command is timed twice.

\begin{description}
  \item[Wall] is measured by the runner and includes one control-plane round
  trip.
  \item[In-box] is measured inside the sandbox by a shell wrapper that
  brackets the command and writes the elapsed nanoseconds to standard error,
  from where the runner recovers it.
\end{description}

Compute, storage and application workloads report in-box durations. Lifecycle
and the execution dimension report wall durations, because for those the round
trip is the thing under study. A sandbox whose shell cannot produce
nanosecond timestamps yields no in-box value, and the framework then falls back
to wall and records which source it used, rather than silently reporting a
zero.

\subsection{Load generation}

Cluster benchmarks generate load in two modes. Closed-loop holds a fixed
concurrency and starts new work as old work finishes. Open-loop sends requests
on a fixed-rate ticker without waiting for the previous one, so a slow provider
shows up as a growing latency tail rather than as a quietly reduced offered
load.

Two client-side confounders are controlled explicitly. The HTTP transport is
configured with a large connection pool, because with Go's default pool a
burst of a thousand creates queues on the client's own connections and the
benchmark then measures the runner. And HTTP 429 is surfaced as a typed error
rather than retried: a throughput ceiling that is really an account quota must
not be published as a throughput ceiling.

\subsection{Teardown}

Every sandbox is destroyed on a context derived with
\texttt{context.WithoutCancel}, so an interrupted or timed-out run still
releases what it allocated. Cleanup failures are counted and reported; a
provider that leaks sandboxes after a burst is billing for them.

\subsection{Reproducibility}

A run writes one JSON file containing the applied configuration, the client
platform and operator-supplied location, every summary, and every raw sample,
so any statistic in the paper can be recomputed from the artefact. Adapter
configuration blocks hold credentials, so only their \emph{key names} are
echoed. Bootstrap confidence intervals are seeded from the sample values
themselves, which makes a reported interval reproducible from the archived
samples alone. Intervals are suppressed below ten samples rather than
computed on data that cannot support them.

The same JSON renders to Markdown for reading and to a \texttt{booktabs}
tabular for this paper, so no number is retyped between the artefact and
Table~\ref{tab:results}.
